% Section 6, part 6.
\subsection{The lattice theorem}
\label{sec:mod-lattice}

\begin{lem}[Counting]\genusstatus{AI-proposed, unchecked}
\label{lem:mod-count}
For every $w\geq 0$, $V_w$ has exactly $\mathrm d(w)$ basis triples $(m,k,b)$ with $b\neq 1$ and exactly $\mathrm e(w)$ with $b=1$.
\end{lem}

\begin{proof}
Fix $m$ with $5m\leq w$ and put $r:=w-5m$. A triple with $b\neq 1$ needs $w_b+4k=r$, $w_b\in\{1,4,3,2\}$ for $b=0,2,3,4$; these meet every class modulo $4$ once, so there is one candidate $b$, which fails only if $w_b>r$, that is $r=0$, $w=5m$. This gives $\mathrm d(w)$. The triples with $b=1$ ($w_1=0$) are the pairs $(m,k)$ with $5m+4k=w$.
\end{proof}

\begin{lem}[The leg class]\genusstatus{AI-proposed, unchecked}
\label{lem:mod-classL}
\begin{enumerate}
\item $\bbQ[p,\bar p]=\bbQ[p,x]=\bbQ[v][p]+\bbQ[v][x]$, and $\bbQ[p,x]$ is a free $\bbQ[v,x]$-module with basis $1,p,p^2,p^3,p^4$.
\item $\bbQ[v][p^{\pm1}]\cap\bbQ(v)[x]=\bbQ[v][x]$.
\item $\cL\cap\bbQ[v][p^{\pm1}]=\bbQ[p,\bar p]$: a Laurent polynomial with coefficients in $\bbQ[v]$ lies in $\cL$ if and only if each of its terms $v^mp^j$ has $j\geq -m$.
\item $\cL/\bbQ(v)[x]=\bbQ[p,\bar p]/\bbQ[v,x]=\bbQ[v][p]/\bbQ[v]=\bigoplus_{i=1}^4\bbQ[v,x]\,p^i$; it is free over $\bbQ[v]$, of dimension $\mathrm d(w)$ in weight $w$.
\end{enumerate}
\end{lem}

\begin{proof}
(i) $\bar p=x-p^4/4$ gives $\bbQ[p,\bar p]=\bbQ[p,x]$, and $v=px-p^5/4$, so $p,x$ are algebraically independent. The map $\bbQ[v,x][T]/(T^5-4xT+4v)\to\bbQ[p,x]$, $T\mapsto p$, is onto, and its source is $\bbQ[x,T]$, so it is an isomorphism, which gives the basis. Clearly $\bbQ[v][p]+\bbQ[v][x]\subset\bbQ[p,\bar p]$. Conversely, if $F\in\bbQ[p,\bar p]$ is homogeneous and $v^mp^{-n}$ ($n>0$, $m\geq n$) is its term with the most negative exponent, unique by homogeneity, subtract a multiple of $v^{m-n}x^n$, whose lowest term is $v^np^{-n}$, and induct. (ii) If $P=\sum_{n\leq N}a_nx^n\in\bbQ[v][p^{\pm1}]$ with $a_N\neq 0$, only $x^N$ contributes to $p^{4N}$, with coefficient $a_N4^{-N}$, so $a_N\in\bbQ[v]$; induct on $N$. (iii) $\supset$ follows from (i). If $F=R+P$, $R\in\bbQ[v][p]$, $P\in\bbQ(v)[x]$, is a Laurent polynomial over $\bbQ[v]$, then $P\in\bbQ[v][x]$ by (ii), and $F\in\bbQ[p,\bar p]$ by (i). (iv) $\bbQ[p,\bar p]\to\cL/\bbQ(v)[x]$ is onto, with kernel $\bbQ[p,\bar p]\cap\bbQ(v)[x]=\bbQ[v,x]$ by (ii); the quotient is $\bigoplus_{i=1}^4\bbQ[v,x]p^i$ by (i). As the lowest term of $x^n$ is $v^np^{-n}$, $\bbQ[v][p]\cap\bbQ[v][x]=\bbQ[v]$, so the quotient is also $\bbQ[v][p]/\bbQ[v]$, with the basis $v^mp^{w-5m}$, $5m<w$, in weight $w$: $\mathrm d(w)$ elements.
\end{proof}

\begin{lem}[Reduction modulo $v$]\genusstatus{AI-proposed, unchecked}
\label{lem:mod-modv}
In every weight, $\Lambda\otimes\bbQ[v]/(v)\colon\mathbb F/v\mathbb F\to\mathsf N/v\mathsf N$ is an isomorphism; here $\Lambda$ is followed by the projection $\mathsf M\to\mathsf N$.
\end{lem}

\begin{proof}
$\Lambda$ is graded, $\bbQ[v]$-linear, with values in $\mathsf M$ (\cref{lem:mod-closure}). \emph{Source.} $(\mathbb F/v\mathbb F)_w$ has the basis $\psi^kH^b$ with $w_b+4k=w$: one element if $4\nmid w$; $\{\psi^{n-1}H^2,\psi^nH\}$ if $w=4n\geq 4$; $\{H\}$ if $w=0$. \emph{Target.} $\mathsf N/v\mathsf N=\bbQ[p,\bar p]/(\bbQ[v,x]+p\bar p\,\bbQ[p,\bar p])\oplus\log p\cdot\bbQ[x]$. Here $\bbQ[p,\bar p]/(p\bar p)$ has the basis $1$, $p^a$, $\bar p^c$ ($a,c\geq 1$), and $\bbQ[v,x]$ maps to the span of $1$ and $x^n\equiv p^{4n}/4^n+\bar p^n$; so $(\mathsf N/v\mathsf N)_w$ has the basis $\{p^w\}$ if $w\geq 1$, $4\nmid w$; $\{p^{4n},\log p\cdot x^n\}$ if $w=4n\geq 4$; $\{\log p\}$ if $w=0$. \emph{The map.} By \cref{lem:mod-closure}~(i), $I(p^a)\equiv -p^{a+4}/(a+4)$ mod $v\mathsf M$ ($a\geq 1$), so $\phi^b_k\equiv(-1)^kp^{w_b+4k}/\prod_{i=1}^k(w_b+4i)$ for $b=0,4,3$; again by \cref{lem:mod-closure}~(i), $I(\bar p^c)\equiv-\bar p^{c+1}/(c+1)$ mod $v\mathsf M$ for $c\geq1$ (the other term is $-vp^3\bar p^{c-1}/(4-c)$, or $-v^4\log p$ if $c=4$), so $\phi^2_k\equiv(-1)^k\bar p^{k+1}/(k+1)!$ mod $v\mathsf M$, which is log-free; and $x^n\equiv p^{4n}/4^n+\bar p^n$ modulo $p\bar p\,\bbQ[p,\bar p]$ gives $\bar p^n\equiv-p^{4n}/4^n$ in $\mathsf N/v\mathsf N$, hence $\phi^2_k\equiv(-1)^{k+1}p^{4k+4}/\big(4\prod_{i=1}^k(4+4i)\big)$. The $\log p$-component of $\phi^1_k$ is $(-1)^kx^k/k!$ modulo $v$, by induction on $k$: for the part $\log p\cdot g(x)$ this is \cref{lem:mod-closure}~(ii) with $G$ without constant term (the terms $v^c\log p$ of $\int G\,\dd p/p$ then have $c\geq1$), and by \cref{lem:mod-closure}~(i) $I$ produces $\log p$ from $\bbQ[p,\bar p]$ only with a factor $v$. In weight $4n$ the matrix is triangular with non-zero diagonal, in the other weights a non-zero $1\times 1$ matrix.
\end{proof}

\begin{thm}[The lattice theorem]\genusstatus{AI-proposed, unchecked}
\label{thm:mod-lattice}
For every $w\geq 0$:
\begin{enumerate}
\item[(a)] $V_w$ has $\mathrm d(w)$ basis triples with $b\neq 1$ and $\mathrm e(w)$ with $b=1$;
\item[(b)] for $X\in\sKo\cap V_w$, $\Lambda(X)$ lies in $\bbQ[p,\bar p]=\bbQ[v][p]+\bbQ[v][x]$, exactly and not only modulo functions of $x$;
\item[(c)] $\Lambda\colon\sKo\cap V_w\to(\cL/\bbQ(v)[x])_w$ is bijective, and both sides have dimension $\mathrm d(w)$; here $\cL/\bbQ(v)[x]$ is graded through its identification with $\bbQ[p,\bar p]/\bbQ[v,x]$ (\cref{lem:mod-classL}~(iv));
\item[(d)] $\Lambda\colon V_w\to\mathsf N_w$ is bijective; $\mathsf N$ is a free $\bbQ[v]$-module with basis $\{\Lambda(\psi^kH^b)\}$.
\end{enumerate}
The coefficients must lie in $\bbQ[v]$: over $\bbQ(v)$ the statement fails, since $\Lambda(H^2/v)=-1/p\notin\cL$.
\end{thm}

\begin{proof}
(a) is \cref{lem:mod-count}. (b) $\Lambda(X)\in\mathsf M$ (\cref{lem:mod-closure}) is log-free (\cref{prop:mod-dii}), so it lies in $\bbQ[p,\bar p]$, and \cref{lem:mod-classL}~(i) applies. (d) Let $C$ be the cokernel of the graded $\bbQ[v]$-linear map $\Lambda\colon\mathbb F\to\mathsf N$, whose pieces are finite-dimensional and zero in negative weights. By right exactness and \cref{lem:mod-modv}, $C/vC=0$; if $w$ were minimal with $C_w\neq 0$, then $(vC)_w=vC_{w-5}=0$ and $C_w=0$. So $\Lambda$ is onto, and as $\dim V_w=\mathrm d(w)+\mathrm e(w)=\dim\mathsf N_w$ (\cref{lem:mod-count}, \cref{lem:mod-classL}~(iv), and $\dim\bbQ[v,x]_w=\mathrm e(w)$), $\Lambda_w$ is bijective. (c) By (b) and \cref{lem:mod-classL}~(iv) the map is well defined. If $\Lambda(X)\in\bbQ(v)[x]$, then $\Lambda(X)\in\bbQ[v,x]$ (\cref{lem:mod-classL}~(ii)), so $\Lambda(X)=0$ in $\mathsf N$ and $X=0$ by (d). Given $\phi\in(\bbQ[p,\bar p]/\bbQ[v,x])_w$, $X:=\Lambda^{-1}(\phi\oplus 0)$ has a log-free leg, so $X\in\sKo$ (\cref{prop:mod-dii}). The dimensions are \cref{lem:mod-classL}~(iv). Finally $\Lambda(H^2)=\varrho_2^{-1}v/p=-v/p$, and $-1/p\notin\cL$ by \cref{lem:mod-classL}~(iii).
\end{proof}

\begin{cor}[Splitting of $V_w$]\genusstatus{AI-proposed, unchecked}
\label{cor:mod-split}
$V_w=(\sKo\cap V_w)\oplus\operatorname{span}\{v^m\psi^kH:5m+4k=w\}$. Every $X\in V_w$ is $X_0+\sum a_{m,k}v^m\psi^kH$ with $X_0\in\sKo$, the $a_{m,k}$ being read off from the $\log p$-coefficient of $\Lambda(X)$.
\end{cor}

\begin{proof}
By \cref{lem:mod-closure}~(ii) and induction, the $\log p$-coefficient of $\phi^1_k$ is $(-1)^kx^k/k!$ plus terms of lower degree in $x$, so $X\mapsto$ the $\log p$-coefficient of $\Lambda(X)$ is triangular, hence bijective, from the span of the $v^m\psi^kH$ onto $\bbQ[v,x]_w$ ($v^m\psi^kH\leftrightarrow v^mx^k$); its kernel on $V_w$ is $\sKo\cap V_w$ by \cref{prop:mod-dii}.
\end{proof}

\begin{cor}[Splitting of $\sK\cap V_w$]\genusstatus{AI-proposed, unchecked}
\label{cor:mod-Ksplit}
$\sK\cap V_w=(\sKo\cap V_w)\oplus\bbQ v^{w/5}H$ if $5\mid w$, and $\sK\cap V_w=\sKo\cap V_w$ otherwise. On $\sK$ the only legs beyond $\bbQ[p,\bar p]$ are multiples of $\log p$: for $X\in\sK\cap V_w$, $\Lambda(X)=\Lambda(X-\Pi_0(X)H)+\Pi_0(X)\log p$ with $\Lambda(X-\Pi_0(X)H)\in\bbQ[p,\bar p]$.
\end{cor}

\begin{proof}
By \cref{prop:mod-dii}, $X\in\sK$ if and only if the $\log p$-coefficient of $\Lambda(X)$ equals $\Pi_0(X)\in\bbQ v^{w/5}$ (\cref{def:tra-tower}; zero unless $5\mid w$). Write $X=X_0+Y$ as in \cref{cor:mod-split}; $\Lambda(Y)$ and $\Lambda(\Pi_0(X)H)=\Pi_0(X)\log p$ have the same $\log p$-coefficient, so $Y=\Pi_0(X)H$ by injectivity on the span of the $v^m\psi^kH$. Conversely $\sKo\subset\sK$ and $H\in\sK$. The last statement is \cref{thm:mod-lattice}~(b) for $X-\Pi_0(X)H\in\sKo\cap V_w$.
\end{proof}

% Section 6, part 7.
\subsection{L-regularity and the Gap Theorem}
\label{sec:mod-Lreg}

We use the closed formula of \cite[Thm.~7.1]{ABDKS25}: for a rational function $x(z)$ on $\bbP^1$ such that $\dd x$ has only simple zeros, all at finite points (the theorem is printed for an arbitrary rational $x$; the condition is that of the preamble of \cite[\S7.1]{ABDKS25} and of \cite[Theorem~1.14]{ABDKS25}, from which the theorem is derived), the differentials of topological recursion for $(x=x(z),y=z)$ with the Bergman kernel and $\omega^{(0)}_{1,0}=-y\,\dd x$ are, for $(g,m)\neq(0,1)$,
\begin{equation}
\label{eq:mod-abdks}
\omega^{(g)}_{m,0}=[\hbar^{2g}]\sum_{r_1,\dots,r_m\geq 0}\prod_{i=1}^m\big(\mathcal D_{x_i}\big)^{r_i}[w_i^{r_i}]\sum_{\Gamma}\hbar^{2g(\Gamma)}\prod_{\{i,j\}\in E(\Gamma)}W_{i,j}\prod_{i=1}^mW_i
\end{equation}
\cite[(7.9)]{ABDKS25}, where $\Gamma$ runs over the connected simple graphs on the vertices $1,\dots,m$, $g(\Gamma)$ is the first Betti number, $\mathcal D_{x_i}\eta:=\dd(\eta/\dd x_i)$ on $1$-forms in $z_i$, $S(u):=\sinh(u/2)/(u/2)$,
\begin{gather*}
W_i:=-\exp\!\big(-w_i(S(w_i\hbar\partial_{z_i})-1)\,x(z_i)\big)\frac{\dd z_i}{w_i},\\
W_{i,j}:=\frac{1}{\hbar^2}\Big(\exp\!\big(\hbar^2w_iw_jS(w_i\hbar\partial_{z_i})S(w_j\hbar\partial_{z_j})(z_i-z_j)^{-2}\big)-1\Big)
\end{gather*}
\cite[(7.6), (7.7)]{ABDKS25}, whose closed form $W_{i,j}=w_iw_j\big/\bigl((z_i-z_j)^2-\hbar^2(w_i+w_j)^2/4\bigr)$, printed in \cite[(7.7)]{ABDKS25}, expands as $W_{i,j}=\sum_{k\geq 0}(\hbar^2/4)^k(w_i+w_j)^{2k}w_iw_j(z_i-z_j)^{-2k-2}$; for one point \cite[(7.10)]{ABDKS25},
\begin{equation}
\label{eq:mod-710}
\sum_{g\geq 1}\hbar^{2g}\omega^{(g)}_{1,0}=\sum_{r\geq 0}\mathcal D_x^r[w^r]W_1 .
\end{equation}
The standing assumptions of \cite[\S1.1]{ABDKS25} (a compact curve with its Bergman kernel; the zeros of $\dd x$ simple; $y$ regular with $\dd y\neq0$ there) and the hypotheses of \cite[Theorem~1.14]{ABDKS25} for the swapped data $(\bbP^1,p,\Wv)$, whose recursion is trivial because $\dd p$ has no zeros, hold for $x=\Wv$, $y=p$, $v\neq 0$ (\cref{def:mod-curve}), and $\omega_{h,n}=(-1)^n\omega^{(h)}_{n,0}$ (\cref{conv:mod-kernel}).

\begin{lem}[Structure of the correlators]\genusstatus{AI-proposed, unchecked}
\label{lem:mod-abdks}
For $2h-2+n>0$ there are an integer $K\geq 0$ and $N_{h,n}\in\bbQ[v][p_1,\dots,p_n]$ with
\[
\omega_{h,n}=N_{h,n}(p;v)\,\frac{\prod_i\dd p_i}{\prod_i(p_i^5-v)^K}.
\]
\end{lem}

\begin{proof}
(1) \emph{For fixed $(h,n)$, \eqref{eq:mod-abdks} is a finite sum of terms in $\bbQ[v][p_i,p_i^{-1},(p_i^5-v)^{-1},(p_i-p_j)^{-1}]\prod_i\dd p_i$.} Since $S(u)-1=\sum_{k\geq 1}u^{2k}/(4^k(2k+1)!)$,
\[
-w\big(S(w\hbar\partial_p)-1\big)\Wv(p)=\sum_{k\geq 1}w^{2k+1}\hbar^{2k}\epsilon_k(p),\qquad \epsilon_k:=-\frac{\Wv^{(2k)}}{4^k(2k+1)!}\quad(\text{local letters}),
\]
with $\Wv^{(2k)}=(p^4/4)^{(2k)}+(2k)!\,v\,p^{-2k-1}\in\bbQ[v][p,p^{-1}]$. In $W_i$ every $\hbar^{2k}$ comes with $w_i^{2k+1}$, in $W_{i,j}$ every $\hbar^2$ with $(w_i+w_j)^2$ and a factor $(p_i-p_j)^{-2}$; so at fixed order in $\hbar$ the degrees in the $w_i$ are bounded, finitely many $r_i$ and graphs contribute, and the coefficients are polynomials in the $\epsilon_k(p_i)$ and the $(p_i-p_j)^{-1}$. Finally $\mathcal D_x(f\,\dd p)=\dd\big(p^2f/(p^5-v)\big)$ only adds powers of $(p^5-v)^{-1}$. (2) \emph{Clearing.} Write $\omega_{h,n}=P\prod_i\dd p_i\big/\big(\prod_ip_i^{a}(p_i^5-v)^{K}\prod_{i<j}(p_i-p_j)^{e}\big)$ with $P\in\bbQ[v][p]$. By \cref{conv:mod-kernel} there are no poles at $p_i=0$ or on the diagonals, so for every $v\neq 0$, $P$ vanishes to order $\geq e$ along $p_i=p_j$ and $\geq a$ along $p_i=0$. Dividing by the monic polynomials $(p_i-p_j)^e$, $p_i^a$ in $p_i$ over $\bbQ[v][p_k:k\neq i]$, the remainder vanishes for all $v\neq 0$ and $p$, hence is $0$; one factor at a time, these denominators cancel in $\bbQ[v][p]$, without negative powers of $v$.
\end{proof}

\begin{thm}[L-regularity]\genusstatus{AI-proposed, unchecked}
\label{thm:mod-Lreg}
Let $2h-2+n>0$. For legs $\phi_1,\dots,\phi_n\in\bbQ[v][p]$ the residue pairing $\langle\phi_1\cdots\phi_n\rangle^{\mathrm{TR}}_h$ lies in $\bbQ[v]$. For legs in $\cL=\bbQ[v][p]+\bbQ(v)[x]$ the residue pairing, defined for $v$ off the poles of the coefficients, is the restriction of an element of $\bbQ[v]$; the components in $\bbQ(v)[x]$ pair to zero.
\end{thm}

\begin{proof}
By \cref{lem:mod-abdks}, $\phi_l(p_l)\omega_{h,n}$ is a rational $1$-form in $p_l$ with poles at the branch points and possibly at $\infty$, holomorphic at $p_l=0$; so $\sum_b\Res_{p_l=\varsigma_b}=-\Res_{p_l=\infty}$. Expanding $(p^5-v)^{-K}=\sum_{j\geq 0}\binom{K+j-1}{j}v^jp^{-5K-5j}$ for $|p|^5>|v|$, $\Res_{p=\infty}p^a\,\dd p/(p^5-v)^K$ is $-\binom{K+j-1}{j}v^j$ if $a+1=5(K+j)$ with an integer $j\geq 0$, and $0$ otherwise, and the result is again of the form of \cref{lem:mod-abdks} in the other variables. The iterated residues commute (no diagonal poles), so
\[
\langle\phi_1\cdots\phi_n\rangle^{\mathrm{TR}}_h=(-1)^n\Res_{p_1=\infty}\cdots\Res_{p_n=\infty}\Big[\prod_l\phi_l(p_l)\,N_{h,n}\prod_l\dd p_l\Big/\prod_l(p_l^5-v)^K\Big]\in\bbQ[v].
\] The components in $\bbQ(v)[x]$ pair to zero by \cref{lem:mod-xfun}.
\end{proof}

\begin{prop}[The period of $\omega_{h,1}$]\genusstatus{AI-proposed, unchecked}
\label{prop:mod-period}
Let $h\geq 1$. Then:
\begin{enumerate}
\item $\omega_{h,1}=\dd\Phi_h$ with $\Phi_h$ rational and holomorphic at $p=0$ and $p=\infty$;
\item $\int_\infty^0\omega_{h,1}=\frac{B_{2h}}{2h}\,v^{1-2h}$ along any path avoiding the branch points;
\item for $h\geq 2$, $\FTR_h=\dfrac{B_{2h}}{2h(2h-2)}\,v^{2-2h}$.
\end{enumerate}
\end{prop}

\begin{proof}
\emph{Step 1 (exactness).} Put $\omega^{\mathrm A}_{h,1}:=\omega^{(h)}_{1,0}=-\omega_{h,1}$. By \eqref{eq:mod-710}, $\omega^{\mathrm A}_{h,1}=[\hbar^{2h}]\sum_{r\geq 1}\mathcal D_x^r[w^r]W_1$; the term $r=0$ is $-[w^1]e^{E}\dd p=0$, as $E:=\sum_kw^{2k+1}\hbar^{2k}\epsilon_k=O(w^3)$. For $r\geq 1$, $\mathcal D_x^r(a\,\dd p)=\dd\big(D_x^{r-1}(a/x')\big)$ (\cref{def:mod-xi}), so $\omega^{\mathrm A}_{h,1}=\dd\Phi^{\mathrm A}_h$ with $\Phi^{\mathrm A}_h:=\sum_{r\geq 1}D_x^{r-1}(a_r/x')$, $a_r:=-[\hbar^{2h}][w^{r+1}]e^{E}=-\sum\frac1{j!}\prod_{i=1}^j\epsilon_{k_i}$, the finite sum over the ordered $j$-tuples of positive integers with $k_1+\dots+k_j=h$, where $j:=r-2h+1$ (the sum is empty unless $1\leq j\leq h$). Put $\Phi_h:=-\Phi^{\mathrm A}_h$.

\emph{Step 2 ($p=\infty$).} $\deg\epsilon_k\leq 4-2k$, $\deg(1/x')=-3$, and $D_x$ lowers the degree by $4$, so every term of $\Phi^{\mathrm A}_h$ has degree at most $(4j-2h)-3-4(2h-2+j)=5-10h<0$: $\Phi^{\mathrm A}_h(\infty)=0$.

\emph{Step 3 ($p=0$).} The \emph{pure part} of a term replaces $\epsilon_k$, $1/x'=-(p^2/v)(1-p^5/v)^{-1}$ and $D_x$ by their leading parts $\epsilon^0_k:=-v\,p^{-2k-1}/(4^k(2k+1))$, $-p^2/v$ and $\delta_0:=-(p^2/v)\,\dd/\dd p$; in $\bar p=v/p$, $\delta_0=\dd/\dd\bar p$, $-p^2/v=-v/\bar p^2$, $\epsilon^0_k=-\bar p^{2k+1}v^{-2k}/(4^k(2k+1))$. With the local letters $\tilde w:=w\bar p$, $\hbar':=\hbar/(2v)$, $E_0:=-\sum_{k\geq 1}\hbar'^{2k}\tilde w^{2k+1}/(2k+1)$ and $C_{h,r}:=[\hbar'^{2h}][\tilde w^{r+1}]\exp(E_0)$, the pure part of $a_r$ is $-(2v)^{-2h}\bar p^{r+1}C_{h,r}$, and that of the $r$-th term is $(\dd/\dd\bar p)^{r-1}\big[(2v)^{-2h}v\,C_{h,r}\bar p^{r-1}\big]=v(2v)^{-2h}(r-1)!\,C_{h,r}$, a constant. Replacing a leading part by the full factor adds terms of order higher in $p$ by at least $3$ (by $\geq 2k+1$ for $\epsilon_k$, $\geq 5$ for $1/x'$ and for $D_x$), and $\delta_0$ never lowers it; so $\Phi^{\mathrm A}_h$ is holomorphic at $0$ with $\Phi^{\mathrm A}_h(0)=v(2v)^{-2h}\Sigma_h$, $\Sigma_h:=\sum_r(r-1)!\,C_{h,r}$.

\emph{Step 4 ($\Sigma_h=-4^hB_{2h}/(2h)$).} Let $\mathsf a(\mathsf x):=\sum_{k\geq 1}\mathsf x^k/(2k+1)$, so $E_0=-\tilde w\,\mathsf a(\hbar'^2\tilde w^2)$ and, with $r+1=2h+j$, $C_{h,r}=\frac{(-1)^j}{j!}[\mathsf x^h]\mathsf a(\mathsf x)^j$. By the binomial series (the terms $j=0$ and $j>h$ have no $\mathsf x^h$),
\[
\Sigma_h=\sum_{j=1}^h(-1)^j\frac{(2h-2+j)!}{j!}[\mathsf x^h]\mathsf a^j=(2h-2)!\,[\mathsf x^h](1+\mathsf a(\mathsf x))^{1-2h}.
\]
With $\mathsf x=\mathsf y^2$, $1+\mathsf a(\mathsf y^2)=\operatorname{artanh}(\mathsf y)/\mathsf y$, so $[\mathsf x^h](1+\mathsf a)^{1-2h}=[\mathsf y^{1}]\operatorname{artanh}(\mathsf y)^{1-2h}$. The Lagrange--B\"urmann formula for $\mathsf y=\tanh\mathsf u$, $[\mathsf y^n]H(\mathsf u(\mathsf y))=\frac1n\Res_{\mathsf u=0}H'(\mathsf u)\tanh(\mathsf u)^{-n}\dd\mathsf u$ ($n\neq 0$, $H$ a Laurent series; change of variables and one integration by parts), with $H=\mathsf u^{1-2h}$, $n=1$, gives $(1-2h)[\mathsf u^{2h}](\mathsf u\coth\mathsf u)=(1-2h)4^hB_{2h}/(2h)!$. So $\Sigma_h=-4^hB_{2h}/(2h)$ and $\Phi^{\mathrm A}_h(0)=-\frac{B_{2h}}{2h}v^{1-2h}$. This proves (i).

\emph{Step 5.} (ii) $\int_\infty^0\omega_{h,1}=\Phi_h(0)-\Phi_h(\infty)=-\Phi^{\mathrm A}_h(0)$. (iii) As $\Phi_{y\dd x}=p^5/5-v\log p$, $(2-2h)\FTR_h=\frac15\langle p^5\rangle^{\mathrm{TR}}_h+v\langle-\log p\rangle^{\mathrm{TR}}_h$ (\cref{def:mod-pair}). Here $\langle p^5\rangle^{\mathrm{TR}}_h\in\bbQ[v]$ (\cref{thm:mod-Lreg}) has exponent $2-2h<0$ (\cref{lem:mod-scale}~(ii)), hence vanishes, and $\langle-\log p\rangle^{\mathrm{TR}}_h=\int_{\mathfrak r}\omega_{h,1}=-\frac{B_{2h}}{2h}v^{1-2h}$ by (ii). So $\FTR_h=\frac{-v}{2-2h}\cdot\frac{B_{2h}}{2h}v^{1-2h}$.
\end{proof}

\begin{thm}[The Gap Theorem]\genusstatus{AI-proposed, unchecked}
\label{thm:mod-gap}
On the model curve, for every $h\geq 2$, $\FTR_h=\dfrac{B_{2h}}{2h(2h-2)}v^{2-2h}$ in the convention of \cref{conv:mod-kernel}; hence $\Fm_h=c_hv^{2-2h}$ with $c_h=(-5)^{1-h}\dfrac{B_{2h}}{2h(2h-2)}=\dfrac{(-1)^{h-1}B_{2h}}{2h(2h-2)\,5^{h-1}}$.
\end{thm}

\begin{proof}
\Cref{prop:mod-period}~(iii) and \cref{cor:mod-trF}.
\end{proof}

\begin{rmk}[A first proof through a formula of Hock]
\label{rmk:mod-hock}
The formula \cite[Cor.~1.2, (1.3)]{Hock25}, whose conventions \cite[(2.1), (1.1)]{Hock25} are those of \cref{conv:mod-kernel}, also gives $\FTR_h$: for $y=p$ only the leading terms at $p=0$ contribute, and Step~4 above applies. This route is not used. \draftnote{The closed forms printed in \cite[\S3.2]{Hock25} are not used; only Cor.~1.2 there is cited.}
\end{rmk}

% Section 6, part 8.
\subsection{Theorem V: logarithmic legs}
\label{sec:mod-V}

Here the model curves form a family in $v\neq 0$; $\partial_v|_x$ is the derivative in $v$ at fixed $x$, computed away from the branch points through the inverse branches $p(x,v)$, and $\partial_v|_p$ the derivative at fixed $p$.

\begin{lem}[Fixed $p$ versus fixed $x$]\genusstatus{AI-proposed, unchecked}
\label{lem:mod-fixedx}
$\partial_vp|_x=-p/(p^5-v)$. For every family of meromorphic forms depending holomorphically on $v$, $\partial_v|_x=\partial_v|_p-\mathrm{Lie}_{\cX}$ in each variable, where $\cX:=\frac{p}{p^5-v}\frac{\dd}{\dd p}$. Moreover $\partial_v(y\,\dd x)|_x=-\dd p/p=\int_{q\in\mathfrak r}\omega_{0,2}(\cdot,q)$.
\end{lem}

\begin{proof}
$\partial_vp|_x=-(\partial_vx|_p)/(\partial_px)=-(1/p)\big/\big((p^5-v)/p^2\big)$. For $\omega=f(p,v)\,\dd p$ the chain rule gives $\partial_v|_x\omega=(\partial_vf+f_p\,\partial_vp|_x)\dd p+f\,\dd(\partial_vp|_x)=\partial_v|_p\omega-\mathrm{Lie}_{\cX}\omega$, in each variable. Since $y=p$, $\partial_v(y\,\dd x)|_x=(\partial_vp|_x)\,\dd x=-\dd p/p$, and $\int_{q\in\mathfrak r}\frac{\dd p\,\dd q}{(p-q)^2}=\dd p\big[\frac{1}{p-q}\big]_{q=0}^{q=\infty}=-\frac{\dd p}{p}$.
\end{proof}

\begin{lem}[Integration by parts]\genusstatus{AI-proposed, unchecked}
\label{lem:mod-ibp}
Let $w$ be a meromorphic $1$-form on $\bbP^1$ with poles only at the branch points, zero residues, holomorphic at $p=0$ and $O(\dd p/p^2)$ at $p=\infty$ (for instance $\omega_{h,n}$ in one slot, $2h-2+n>0$).
\begin{enumerate}
\item For $\phi\in\bbQ[v][p]$ let $(\cX\dd\phi)_+\in\bbQ[v][p]$ be the quotient in the division $p\,\phi'(p)=(\cX\dd\phi)_+\,(p^5-v)+r$, $\deg r<5$; for $\phi=p^m$, $(\cX\dd\phi)_+=m\sum_{j\geq 0,\ m-5-5j\geq 0}v^jp^{m-5-5j}$. Then, for $R>|\varsigma|$, $(2\pi\ri)^{-1}\oint_{|p|=R}\phi\,\mathrm{Lie}_{\cX}w$ equals $-\langle(\cX\dd\phi)_+,w\rangle$.
\item $\int_{\mathfrak r}\mathrm{Lie}_{\cX}w=0$.
\end{enumerate}
\end{lem}

\begin{proof}
(i) $\mathrm{Lie}_{\cX}(\phi w)=\dd(\iota_{\cX}(\phi w))$ has zero circle integral, and $\mathrm{Lie}_{\cX}(\phi w)=(\cX\phi)\,w+\phi\,\mathrm{Lie}_{\cX}w$ with $\cX\phi=(\cX\dd\phi)_++r/(p^5-v)$. The form $\frac{r}{p^5-v}w$ has all its poles inside the circle and is $O(\dd p/p^3)$ at $\infty$, so its circle integral vanishes; and the circle pairing of a polynomial is its residue pairing (\cref{def:mod-pair}). For $\phi=p^m$, expand $mp^m/(p^5-v)$ for $|p|^5>|v|$. (ii) $\int_{\mathfrak r}\mathrm{Lie}_{\cX}w=[\iota_{\cX}w]_0^\infty=0$: $\cX$ vanishes at $0$, where $w$ is holomorphic, and $\cX=O(p^{-4})\,\dd/\dd p$, $w=O(\dd p/p^2)$ at $\infty$.
\end{proof}

\begin{thm}[The variational formula, in the form used]\genusstatus{AI-proposed, unchecked}
\label{thm:mod-var}
Let $J=(p_1,\dots,p_n)$ be points off the branch points.
\begin{enumerate}
\item For $n\geq 1$ and $h+n>1$: $\partial_v|_x\,\omega_{h,n}(J)=\int_{q\in\mathfrak r}\omega_{h,n+1}(J,q)$.
\item (Not used.) For $h\geq 2$: $\partial_v\FTR_h=\int_{\mathfrak r}\omega_{h,1}$.
\end{enumerate}
\end{thm}

\begin{proof}
(i) This is \cite[Thm.~5.1]{Eynard:2007kz}, (5-14) for $g+k>1$, for the temperature variation of \cite[\S5.3, (5-27)--(5-30)]{Eynard:2007kz}, translated by \cref{conv:mod-kernel}. (The path of \cite[(5-27)]{Eynard:2007kz}, printed from $\alpha$ to $\alpha'$ with $\Lambda=1$ (letters of that paper), is taken with $\alpha=\infty$, $\alpha'=0$, that is, along $\mathfrak r$ from $0$ to $\infty$ with $\Lambda=-1$.) There a variation at fixed $x$ is $\delta_\Omega y|_x\,\dd x=-\Omega$ \cite[(5-1)]{Eynard:2007kz} with $\Omega(p)=\int_{\partial\Omega}B(p,q)\Lambda(q)$ \cite[(5-13)]{Eynard:2007kz}, and (5-14) reads $\delta_\Omega W^{(h)}_n(J)|_{x}=\int_{\partial\Omega}W^{(h)}_{n+1}(J,q)\Lambda(q)$; on $\bbP^1$ there is no period matrix and no parameter $\kappa$, so the covariant variation $D_\Omega$ of \cite[Thm.~5.1]{Eynard:2007kz} is $\delta_\Omega$. By \cref{lem:mod-fixedx}, $\Omega=-\partial_v(y\,\dd x)|_x=-\int_{\mathfrak r}\omega_{0,2}(\cdot,q)$: the path $\partial\Omega=\mathfrak r$ joins the poles $0$, $\infty$ of $y\,\dd x$, with $\Lambda=-1$. With $\omega_{h,n}=(-1)^nW^{(h)}_n$ this gives $\partial_v|_x\omega_{h,n}=(-1)^{n}(-1)^{n+1}(-1)\int_{\mathfrak r}\omega_{h,n+1}=\int_{\mathfrak r}\omega_{h,n+1}$. The hypotheses hold: the branch points are simple with $\dd y\neq 0$; the variation is at fixed $x$ with the Bergman kernel of the fixed $\bbP^1$; $\Omega/\dd y=1/p$ is regular at the branch points \cite[the assumption before (5-4)]{Eynard:2007kz}; the path avoids them, without hard edges \cite[(5-13), footnote~8]{Eynard:2007kz}; the integrals converge ($\omega_{h,n+1}$ is holomorphic at $0$, $\infty$); and residues at the moving branch points commute with $\partial_v|_x$ (fixed circles in the $x$-plane).

(ii) This part is not used. By \cref{prop:mod-period}~(iii), $\partial_v\FTR_h=-\frac{B_{2h}}{2h}v^{1-2h}$, and by \cref{prop:mod-period}~(ii), $\int_{\mathfrak r}\omega_{h,1}=-\int_\infty^0\omega_{h,1}=-\frac{B_{2h}}{2h}v^{1-2h}$. (It is also the translation of \cite[(5-15)]{Eynard:2007kz}, which is not used.)
\end{proof}

\begin{lem}[Genus zero]\genusstatus{AI-proposed, unchecked}
\label{lem:mod-genus0}
Put $e_{b,k}(p):=\dd p/(p-\varsigma_b)^{k+1}$ and $r_b:=1/(5\varsigma_b^3)$. Then:
\begin{enumerate}
\item $\partial_v|_x\,\omega_{0,2}(p_1,p_2)=\sum_br_b\,e_{b,1}(p_1)\,e_{b,1}(p_2)$;
\item $\omega_{0,3}(p_0,p_1,p_2)=-\sum_be_{b,1}(p_0)e_{b,1}(p_1)e_{b,1}(p_2)/x''(\varsigma_b)$, the negative of \cite[Thm.~4.1 (4-20)]{Eynard:2007kz};
\item $\int_{q\in\mathfrak r}\omega_{0,3}(p_1,p_2,q)=\partial_v|_x\,\omega_{0,2}(p_1,p_2)$;
\item for $\phi_1,\phi_2,\phi_3\in\bbQ[v][p]$,
\[
\begin{aligned}
\langle\phi_1\phi_2\phi_3\rangle^{\mathrm{TR}}_0&=-\tfrac15\textstyle\sum_b\varsigma_b^{-2}\prod_i\phi_i'(\varsigma_b), &\langle\phi_1\phi_2,-\log p\rangle^{\mathrm{TR}}_0&=\tfrac15\textstyle\sum_b\varsigma_b^{-3}\phi_1'(\varsigma_b)\phi_2'(\varsigma_b),\\
\langle\phi_1,(-\log p)^2\rangle^{\mathrm{TR}}_0&=-\tfrac15\textstyle\sum_b\varsigma_b^{-4}\phi_1'(\varsigma_b), &\langle(-\log p)^3\rangle^{\mathrm{TR}}_0&=1/v,
\end{aligned}
\]
and every pairing of $\omega_{0,3}$ with at least one leg in $\bbQ[v][p]$ lies in $\bbQ[v]$.
\end{enumerate}
\end{lem}

\begin{proof}
(i) $\omega_{0,2}$ does not depend on $v$ at fixed $p$, so by \cref{lem:mod-fixedx} $\partial_v|_x\omega_{0,2}=-(\mathrm{Lie}_{\cX,1}+\mathrm{Lie}_{\cX,2})\,\omega_{0,2}$. For a function $g$ and the field $g\,\dd/\dd p$ in both variables,
\[
(\mathrm{Lie}_{1}+\mathrm{Lie}_{2})\frac{\dd p_1\dd p_2}{(p_1-p_2)^2}=\Big(\frac{g'(p_1)+g'(p_2)}{(p_1-p_2)^2}-\frac{2(g(p_1)-g(p_2))}{(p_1-p_2)^3}\Big)\dd p_1\dd p_2=\dd_1\dd_2\frac{g(p_1)-g(p_2)}{p_1-p_2},
\]
and with $g=p/(p^5-v)=\sum_br_b/(p-\varsigma_b)$, $(g(p_1)-g(p_2))/(p_1-p_2)=-\sum_br_b/((p_1-\varsigma_b)(p_2-\varsigma_b))$.
(ii) Only the two terms of ${\sum}'$ with $I_1=\{p_1\}$, $\{p_2\}$ occur. With $t:=q-\varsigma_b$, $\iota_bq=\varsigma_b-t+O(t^2)$, $\dd x(q)=x''(\varsigma_b)t(1+O(t))\dd t$, $\int_{\iota_bq}^q\omega_{0,2}(p_0,\cdot)=2t\,e_{b,1}(p_0)(1+O(t))$ and $y(q)-y(\iota_bq)=2t(1+O(t))$, so $K_b(p_0,q)=\frac{e_{b,1}(p_0)}{2x''(\varsigma_b)t\,\dd t}(1+O(t))$ and the bracket is $-2\,e_{b,1}(p_1)e_{b,1}(p_2)\,\dd t^2(1+O(t))$; the residue is the claim.
(iii) $x''(\varsigma_b)=5\varsigma_b^2$ and $\int_{\mathfrak r}e_{b,1}=\big[-1/(p-\varsigma_b)\big]_0^\infty=-1/\varsigma_b$, so (ii) gives $\sum_br_b\,e_{b,1}(p_1)e_{b,1}(p_2)$, which is (i).
(iv) $\Res_{\varsigma_b}\phi\,e_{b,1}=\phi'(\varsigma_b)$ and $\langle-\log p,e_{b,1}\rangle=-1/\varsigma_b$; inserting into (ii) gives the four formulas ($\frac15\sum_b\varsigma_b^{-5}=1/v$). Constants pair to zero, so with an L-leg we may take $\phi_i\in p\,\bbQ[v][p]$; then only powers $\varsigma_b^N$, $N\geq -4$, occur, and $\sum_b\varsigma_b^N=5v^{N/5}$ if $5\mid N$, $0$ otherwise.
\end{proof}

\begin{prop}[The recursion (V1)]\genusstatus{AI-proposed, unchecked}
\label{prop:mod-V1}
Let $h\geq 0$, $a,k\geq 0$, $n=a+k\geq 1$, $2h-2+n>0$, and $\phi_1,\dots,\phi_a\in\bbQ[v][p]$. Then
\begin{equation}
\label{eq:mod-V1}
\begin{aligned}
\langle\phi_1\cdots\phi_a,(-\log p)^{k+1}\rangle^{\mathrm{TR}}_h={}&\partial_v\langle\phi_1\cdots\phi_a,(-\log p)^k\rangle^{\mathrm{TR}}_h-\sum_{i=1}^a\langle\cdots,\partial_v\phi_i,\cdots,(-\log p)^k\rangle^{\mathrm{TR}}_h\\
&+\sum_{i=1}^a\langle\cdots,(\cX\dd\phi_i)_+,\cdots,(-\log p)^k\rangle^{\mathrm{TR}}_h ,
\end{aligned}
\end{equation}
where $\partial_v\phi_i$ is the derivative of the coefficients and the dots stand for the other legs.
\end{prop}

\begin{proof}
Fix $v_0\neq 0$; for $v$ near $v_0$ take $R>\max_b|\varsigma_b(v)|$ and a ray $\mathfrak r$ avoiding all $\varsigma_b(v)$. Then $C(v):=\langle\phi_1\cdots\phi_a,(-\log p)^k\rangle^{\mathrm{TR}}_h=(2\pi\ri)^{-a}\oint_{|p_i|=R}\int_{\mathfrak r^k}\prod_i\phi_i(p_i)\,\omega_{h,n}$ with contours independent of $v$, and the integrand is holomorphic in $v$, bounded on the rays by a constant times $\prod_j|\dd q_j|/(1+|q_j|)^2$ uniformly near $v_0$: by \cref{lem:mod-abdks}, $\omega_{h,n}=N\prod_i\dd p_i/\prod_i(p_i^5-v)^K$ with $N\in\bbQ[v][p_1,\dots,p_n]$, and holomorphy at $p_i=\infty$ for every $v\neq0$ gives $\deg_{p_i}N\leq5K-2$. The same bound justifies the exchanges of integrals (Fubini) below. Differentiating under the integrals, $\partial_vC=\sum_i\langle\cdots,\partial_v\phi_i,\cdots\rangle^{\mathrm{TR}}_h$ plus the integral of $\prod_i\phi_i\,\partial_v|_p\omega_{h,n}$. By \cref{lem:mod-fixedx} and \cref{thm:mod-var}~(i), $\partial_v|_p\omega_{h,n}=\int_{q\in\mathfrak r}\omega_{h,n+1}(\cdot,q)+\sum_{l=1}^n\mathrm{Lie}_{\cX,l}\,\omega_{h,n}$. The first term gives the pairing with one more log leg (Fubini); a Lie term gives $-\langle\cdots,(\cX\dd\phi_i)_+,\cdots\rangle^{\mathrm{TR}}_h$ in an L-slot and $0$ in a log slot (\cref{lem:mod-ibp}). Rearranging gives \eqref{eq:mod-V1}.
\end{proof}

% Section 6, part 9.
Let $c$ be a branch point, $t:=p-c$, and $\chi$ (a local letter) a local coordinate at $c$ with $x=x(c)+\lambda\chi^2$, $\lambda\neq 0$, so that $\iota_c\colon\chi\mapsto-\chi$. Write $y=\sum_ky_k\chi^k$ and, near $(c,c)$, $\omega_{0,2}=\big[(\chi_1-\chi_2)^{-2}+\mathrm{reg}(\chi_1,\chi_2)\big]\dd\chi_1\dd\chi_2$ with $\mathrm{reg}$ holomorphic and symmetric; then $\mathrm{reg}_{00}:=\mathrm{reg}(0,0)=\frac16\{t;\chi\}(0)$, where $\{t;\chi\}=t'''/t'-\frac32(t''/t')^2$ is the Schwarzian derivative of $t$ as a function of $\chi$ (both sides are $t_3/t_1-(t_2/t_1)^2$ for $t=t_1\chi+t_2\chi^2+t_3\chi^3+\cdots$, by expanding $t'(\chi_1)t'(\chi_2)/(t(\chi_1)-t(\chi_2))^2$). For a $1$-form $\eta=\sum_m\eta_m\chi^m\dd\chi$ holomorphic at $c$ put
\begin{equation}
\label{eq:mod-loc}
\mathrm{Loc}_c[\eta]:=-\frac{1}{4\lambda y_1}\Big[\mathrm{reg}_{00}\,\eta_0+\frac{\eta_2}{12}-\frac{y_3}{4y_1}\,\eta_0\Big].
\end{equation}

\begin{lem}[A local formula for $\omega_{1,1}$]\genusstatus{AI-proposed, unchecked}
\label{lem:mod-w11}
\begin{enumerate}
\item $\omega_{1,1}(p_0)=\sum_c\mathrm{Loc}_c[\omega_{0,2}(p_0,\cdot)]$ for $p_0$ off the branch points, $\omega_{0,2}(p_0,\cdot)$ being a form in the second slot.
\item $\omega_{1,1}=\dd G$ with $G$ rational, with poles only at the branch points and $G(\infty)=0$; so $\omega_{1,1}$ has zero residues, is holomorphic at $p=0$ and is $O(\dd p/p^2)$ at $\infty$.
\item $\langle-\log p\rangle^{\mathrm{TR}}_1=\int_{\mathfrak r}\omega_{1,1}=\sum_c\mathrm{Loc}_c[-\dd q/q]$.
\end{enumerate}
\end{lem}

\begin{proof}
(i) By \eqref{eq:mod-kernel},
\[
\omega_{1,1}(p_0)=\sum_c\Res_{q\to c}K_c(p_0,q)\,\omega_{0,2}(q,\iota_cq).
\]
Expand $\omega_{0,2}(p_0,q')=\dd p_0\sum_mB_m(p_0)\chi^m\dd\chi$ for $q'$ near $c$ with coordinate $\chi$. Then $\int_{\iota_cq}^q\omega_{0,2}(p_0,\cdot)=2\,\dd p_0\sum_{m\text{ even}}B_m\chi^{m+1}/(m+1)$, $y(q)-y(\iota_cq)=2\sum_{k\text{ odd}}y_k\chi^k$ and $\dd x=2\lambda\chi\,\dd\chi$, so
\[
K_c(p_0,q)=\frac{\dd p_0}{\dd\chi}\,\frac{1}{4\lambda y_1}\Big[B_0\chi^{-1}+\Big(\frac{B_2}{3}-B_0\frac{y_3}{y_1}\Big)\chi+O(\chi^3)\Big],
\]
and $\omega_{0,2}(q,\iota_cq)=-\big[(4\chi^2)^{-1}+\mathrm{reg}_{00}+O(\chi^2)\big]\dd\chi^2$ ($\dd(-\chi)=-\dd\chi$, and $\mathrm{reg}(\chi,-\chi)=\mathrm{reg}_{00}+O(\chi^2)$ by symmetry). The coefficient of $\chi^{-1}\dd\chi$ in the product is $-\frac{1}{4\lambda y_1}\big[B_0\,\mathrm{reg}_{00}+\frac14\big(\frac{B_2}{3}-B_0\frac{y_3}{y_1}\big)\big]\dd p_0=\mathrm{Loc}_c[\omega_{0,2}(p_0,\cdot)]$.
(ii) $\omega_{0,2}(p_0,q)=\dd_{p_0}\big[\dd q/(q-p_0)\big]$, and near $c$, $\dd q/(q-p_0)=\sum_mg_m(p_0)\chi^m\dd\chi$ with $g_m$ rational, with poles only at $c$ and $g_m(\infty)=0$. So $B_m\dd p_0=\dd g_m$, and by (i) $\omega_{1,1}=\dd G$ with $G$ a combination of the $g_0,g_2$ at the five branch points. (iii) The residues in (i) are integrals over small circles avoiding $\mathfrak r$; the ray integral converges by (ii), and as $\omega_{0,2}(p_0,q)=O(p_0^{-2})\dd p_0$ uniformly for $q$ on the circles, it may be taken inside: $\int_{p_0\in\mathfrak r}\omega_{0,2}(p_0,q)=-\dd q/q$, and $\mathrm{Loc}_c$ is linear in finitely many Taylor coefficients. Finally, by (i) and \eqref{eq:mod-kernel}, $\mathrm{Loc}_c[\eta]=\Res_{q\to c}\frac12\frac{\int_{\iota_cq}^q\eta}{(y(q)-y(\iota_cq))\dd x(q)}\,\omega_{0,2}(q,\iota_cq)$ for $\eta$ holomorphic at $c$; this form does not depend on the coordinate $\chi$.
\end{proof}

\begin{thm}[Theorem V]\genusstatus{AI-proposed, unchecked}
\label{thm:mod-V}
Let $h\geq 0$, $a,k\geq 0$, $n=a+k\geq 1$, $2h-2+n>0$, and $\phi_1,\dots,\phi_a\in\bbQ[v][p]$.
\begin{enumerate}
\item If $a\geq 1$, then $\langle\phi_1\cdots\phi_a,(-\log p)^k\rangle^{\mathrm{TR}}_h\in\bbQ[v]$.
\item If $a=0$, then $\langle(-\log p)^k\rangle^{\mathrm{TR}}_h=(\FTR_h)^{(k)}(v)$, with $\FTR_h$ of \cref{thm:mod-gap} for $h\geq 2$, $\FTR_1:=-\frac1{12}\log v$ (here $k\geq 1$) and $\FTR_0:=\frac12v^2\log v$ (here $k\geq 3$).
\item The recursion \eqref{eq:mod-V1} holds.
\end{enumerate}
\end{thm}

\begin{proof}
(iii) is \cref{prop:mod-V1}. (i) Fix $h$ and $a\geq 1$ and induct on $k$. The first admissible case is $k=0$ if $2h-2+a>0$ (\cref{thm:mod-Lreg}), and otherwise $\langle\phi_1,(-\log p)^2\rangle^{\mathrm{TR}}_0$ or $\langle\phi_1\phi_2,-\log p\rangle^{\mathrm{TR}}_0$ (\cref{lem:mod-genus0}~(iv)). In the step, the right side of \eqref{eq:mod-V1} consists of a $v$-derivative of, and pairings with, $a$ L-legs in $\bbQ[v][p]$ ($p^5-v$ is monic) and $k$ log legs; all lie in $\bbQ[v]$ by induction.

(ii) For $a=0$, \eqref{eq:mod-V1} reads $\langle(-\log p)^{k+1}\rangle^{\mathrm{TR}}_h=\partial_v\langle(-\log p)^k\rangle^{\mathrm{TR}}_h$, so it suffices to treat the first admissible $k$. \emph{$h\geq 2$, $k=1$.} As in the proof of \cref{prop:mod-period}~(iii), $(2-2h)\FTR_h=v\langle-\log p\rangle^{\mathrm{TR}}_h$, and $\FTR_h$ is a multiple of $v^{2-2h}$ (\cref{thm:mod-gap}), so $\langle-\log p\rangle^{\mathrm{TR}}_h=(\FTR_h)'(v)$. \emph{$h=0$, $k=3$.} \Cref{lem:mod-genus0}~(iv). \emph{$h=1$, $k=1$.} We show $\langle-\log p\rangle^{\mathrm{TR}}_1=-1/(12v)$.

\emph{Reduction to $v=1$.} By \cref{lem:mod-scale}~(ii) with $k=\varsigma$, $m_\varsigma^*\omega^{(v)}_{1,1}=v^{-1}\omega^{(1)}_{1,1}$, and rays go to rays, so $\langle-\log p\rangle^{\mathrm{TR}}_1(v)=v^{-1}\langle-\log p\rangle^{\mathrm{TR}}_1(1)$. Let $v=1$, $\varsigma=1$, $\varsigma_b=\zeta^{-b}$.

\emph{Five equal contributions.} $\rho_b(p):=\zeta^{-b}p$ sends $1$ to $\varsigma_b$, with $x\circ\rho_b=\zeta^bx$, $y\circ\rho_b=\zeta^{-b}y$, $\rho_b^*\omega_{0,2}=\omega_{0,2}$, $\rho_b^*(\dd q/q)=\dd q/q$; hence $\iota_b\circ\rho_b=\rho_b\circ\iota_0$ and $(y(q)-y(\iota q))\,\dd x(q)$ is unchanged. By the coordinate-free form of $\mathrm{Loc}$, $\mathrm{Loc}_{\varsigma_b}[-\dd q/q]=\mathrm{Loc}_1[-\dd q/q]$.

\emph{The contribution at $c=1$.} With $t=p-1$, $x(1+t)=\frac54+\frac52t^2+\frac54t^4-t^5+O(t^6)$, so $\lambda=5/2$, $\chi=t+\frac{t^3}4-\frac{t^4}5+O(t^5)$ and $t=\chi-\frac{\chi^3}4+\frac{\chi^4}5+O(\chi^5)$. Hence $y_1=1$, $y_2=0$, $y_3=-\frac14$; $\{t;\chi\}(0)=-\frac32$ ($t'''(0)=-\frac32$, $t''(0)=0$), so $\mathrm{reg}_{00}=-\frac14$; and $-\dd q/q=-\big(1-\chi+\frac14\chi^2+O(\chi^3)\big)\dd\chi$, so $\eta_0=-1$, $\eta_2=-\frac14$. By \eqref{eq:mod-loc}, $\mathrm{Loc}_1[-\dd q/q]=-\frac1{10}\big[\frac14-\frac1{48}-\frac1{16}\big]=-\frac1{60}$. By \cref{lem:mod-w11}~(iii) the total at $v=1$ is $5\cdot(-\frac1{60})=-\frac1{12}$. So $\langle-\log p\rangle^{\mathrm{TR}}_1=-1/(12v)=(\FTR_1)'(v)$.
\end{proof}

\input{sections/s06_model_10}



